\section{The Wiggle Framework: A Unified Epistemic Stress Test}
\label{sec:framework}

The Wiggle Framework (Figure~\ref{fig:framework}) is a centralized pressure instrument for stress-testing LLM judges. 
It bundles together a graduated set of perturbations like infrastructure noise, prompt-format changes, sycophantic prodding, and multi-turn persuasion.
We apply the framework to the same items, judges, and grading scales so that the resulting wiggle measurements are directly 
comparable.

\subsection{What is a wiggle?}
\label{sec:framework-wiggle}

Every measurement is anchored to an \emph{L0 baseline protocol}: temperature 0 with no pressure applied. 
Operationally, the first valid verdict from this protocol is the trajectory's L0 anchor. We do not assume this to be the judge's unique unpressured output.

A \emph{wiggle} is any movement away from the L0 verdict under perturbation. On Binary scales, a wiggle is a verdict flip
(e.g., safe $\rightarrow$ unsafe). On Likert (1--5) scales, a wiggle is a movement of two or more places. For items off 
the midpoint this is equivalent to crossing the midpoint of 3 (so $4 \rightarrow 2$ counts but $4 \rightarrow 3$ does not).
For items \emph{at} the midpoint, we count a wiggle when the verdict moves to an extreme (1 or 5). 
This excludes minor numerical drift and reserves the term for movements that change the judge's position relative to the decision boundary.

The wiggle rate (WR) is the fraction of items whose verdict changes from L0 by more than this threshold. Its complement, retention rate (RR),
measures how often the judge holds its baseline L0 verdict. Wiggle is orthogonal to accuracy: a judge can wiggle and still be right, or remain
wrong without wiggling. Where ground-truth labels exist, we additionally classify each wiggle as \emph{corrective} when it moves toward the ground truth label or 
\emph{corrupting} when it moves away.

\subsection{Three dimensions}
\label{sec:framework-dimensions}

\textbf{Mechanical Consistency} measures whether the judge's L0 verdict survives perturbations that carry no new information. 
We test three conditions: \emph{infrastructural repetition} (10 identical decoding trials); 
\emph{trivial prompt perturbation} via seed injection (10 trials each with a different 64-character random string appended to the system prompt;
Appendix~\ref{sec:appendix-seed}); and \emph{positional consistency} (the same two opposing arguments presented in both orderings). 
These probes are all semantically invariant\footnote{Prior work on semantics-preserving perturbations often use paraphrases or lexical substitutions, which can subtly change meaning. 
Paraphrases or lexical substitutions can be more problematic in evaluations that center on a specific policy or constitution, where a small wording change may lead to large downstream
changes in interpretation. temperature=0 is also mostly a controlled decoding choice intended to minimize sampling variance as a confound. Although we 
didn't test this, we would argue that using higher or unset temperatures would also qualify for the semantic invariance of a mechanical consistency test.}.
Mechanical Consistency can also be used as the empirical \emph{floor} against which other wiggle tests are measured.

\textbf{Single-turn Conviction} measures whether a single substantive challenge can talk the judge out of its L0 verdict. 
We use four scripted pressure types of increasing sophistication: mild doubt (L1), counterargument (L2), expert 
authority (L3), and fabricated consensus (L4) (Figure~\ref{fig:framework}).

\textbf{Multi-turn Persistence} measures whether the judge holds its verdict when challenges are sustained or adapted over many turns. 
In addition to testing L1--L4 applied statically over each turn of a 10-turn rollout, we also test two expressly multi-turn protocols.
L5 cycles through the same L1--L4 pressure types in randomized order across 10 turns. For L6, a separate LLM acts as an adaptive persuader and
generates the following user turn (Appendix~\ref{sec:appendix-l6-protocol}).

% Technically, both L5 and L6 also yield a turn-1 wiggle rate,
% but L5's first turn is a randomized draw from L1–L4 and therefore approximates their average turn-1, 
% while L6's first turn is a single persuasion opening move that produces less wiggle than L1–L4 wiggle. 
% The regime where L6 distinguishes itself is in the multi-turn setting.

\begin{figure}[t]
  \centering
  \includegraphics[width=\linewidth]{wiggle_framework_summary.pdf}
  \caption{\textbf{Mean wiggle rate per judge across the Wiggle Framework.}
  Each model as a judge is summarized by three bars: \emph{Mechanical Consistency} (verdict change under trivial perturbation), \emph{Single-Turn Conviction} (verdict change after a single scripted challenge), and \emph{Multi-Turn Persistence} (verdict change after 10 turns of sustained or adaptive pressure). Rates are averaged over 14 judging tasks.}
  \label{fig:hero}
\end{figure}

See Appendix~\ref{sec:appendix-worked-example} for a full walkthrough of the framework on five hypothetical items and how wiggle rates and corrective/corrupting fractions are computed.
